\documentclass[12pt]{article}
\usepackage{verbatim}
\usepackage[dvips]{epsfig}
\usepackage{color}
\usepackage{url}
\usepackage[colorlinks=true]{hyperref}

\begin{document}

\section*{Neurospaces: Documentation}

{\bf Related Documentation:}
% start: userdocs-tag-replace-items related-do-nothing
% end: userdocs-tag-replace-items related-do-nothing

\section*{New Document}

The three above headings are included in all documentation by default.

\begin{itemize}
\item[]  {\bf The Documentation Identifier:} This may be replaced by a descriptor of the type of documentation you have developed, for example {\bf Introduction} or {\bf Tutorial}, etc.
\item[]{\bf Related Documentation Heading:} Followed by a hyperlink list of any related documentation.
\item[] {\bf New Document Heading:} Should be replaced with the title of your document. {\bf Note:} In your document (the \LaTeX\,\,source documentation file--{\it NewDocument.tex}) all text below the {\bf New Document} heading should be replaced with your own content. See \href{../document-create/document-create.tex}{\bf Create\,Document} for further information.
\end{itemize}

\subsection*{Help with \LaTeX\,Documentation}

\subsubsection*{Text color}

Text in your document can be colorized  with the command
\begin{verbatim}
    \textcolor{red}{text}
\end{verbatim}
for example \textcolor{red}{\bf This is IMPORTANT}, \textcolor{green}{green} and \textcolor{blue}{blue} are also recognized arguments.

\subsubsection*{Hyperlinks}

Two classes of hyperlink exist in the Neurospaces Documentation System:

\begin{itemize}
\item[{\bf A.}]{\bf Static Hyperlinks:} 
   \begin{itemize}
      \item{\bf Local links:} These are produced with the command:
\begin{verbatim}
   \href{../NewDocument/NewDocument.tex}{NewDocument}}
\end{verbatim}
for example \href{../NewDocument/NewDocument.tex}{NewDocument} references this document.
Note that these links may not work in your \LaTeX\,{\tt pdf} file viewer, but should be testable with the \href{http://get.adobe.com/reader/}{Acrobat Reader} or \href{http://www.adobe.com/products/acrobatpro/tryout.html}{Acrobat Professional} packages.
      \item {\bf Remote links:} These use the {\tt href} command in a slightly different way:
\begin{verbatim}
    \href{http://www.neurospaces.org/}{Neurospaces Website}
\end{verbatim}
for example see \href{http://www.neurospaces.org/}{Neurospaces Documentation}. Note that these links do not wrap automatically at the end of a line even outside of the {\tt verbatim} environment.
\end{itemize}

\item[{\bf B.}]{\bf Reference Hyperlinks:} See \href{../document-create/document-create.tex}{Document Creation} for further details on managing reference hyperlinks.

\end{itemize}

\subsubsection*{Figures}

Figures can easily be added with the following code:

\begin{verbatim}
\begin{figure}[h]
   \centering
   \includegraphics[scale=0.6]{figures/dummyfig.eps}
   \caption{{\bf A Dummy Figure:} Example of Latex code to incorporate
      a figure into documentation.}
   \label{fig:df-1}
\end{figure}
\end{verbatim}

\begin{figure}[h]
  \centering
   \includegraphics[scale=0.5]{figures/dummyfig.eps}
\caption{{\bf A Dummy Figure:} Example of \LaTeX\,\,\,code to incorporate a figure into documentation.}
  \label{fig:df-2}
\end{figure}

\bibliographystyle{plain}
\bibliography{../tex/bib/g3-refs.bib}

\end{document}
